\documentclass[12pt,a4paper]{article}
\usepackage[utf8]{inputenc}
\usepackage[vietnamese]{babel}
\usepackage[left=2cm,right=2cm,top=2cm,bottom=2cm]{geometry}
\usepackage{mathptmx} % Font Times New Roman
\usepackage{enumitem}

\pagestyle{plain}
\linespread{1.15}

\begin{document}

% --- HEADER QUỐC HIỆU & CƠ QUAN BAN HÀNH ---
\noindent
\begin{minipage}[t]{0.45\textwidth}
    \begin{center}
        \small
        CÔNG AN THÀNH PHỐ HÀ NỘI\\
        \textbf{\underline{PHÒNG CẢNH SÁT HÌNH SỰ}}\\[0.1cm]
        Số: 06/BB-HX\\
        \textit{V/v lấy lời khai nhân chứng hàng xóm}
    \end{center}
\end{minipage}
\hfill
\begin{minipage}[t]{0.52\textwidth}
    \begin{center}
        \small
        \textbf{CỘNG HÒA XÃ HỘI CHỦ NGHĨA VIỆT NAM}\\
        \textbf{\underline{Độc lập - Tự do - Hạnh phúc}}\\[0.1cm]
        \textit{Hà Nội, ngày 25 tháng 07 năm 2026}
    \end{center}
\end{minipage}

\vspace{0.6cm}

% --- TIÊU ĐỀ VĂN BẢN ---
\begin{center}
    \textbf{\large BIÊN BẢN LẤY LỜI KHAI NHÂN CHỨNG}\\
    \textit{(Ghi nhận lời khai của các hộ dân lân cận số 14 Đường Bờ Sông)}
\end{center}

\vspace{0.3cm}

\noindent Vào hồi 09 giờ 00 phút đến 11 giờ 00 phút, ngày 25 tháng 07 năm 2026.\\
Cán bộ điều tra: Trung úy Nguyễn Văn Hoàng tiến hành lấy lời khai các nhân chứng:

\section*{I. LỜI KHAI BÀ NGUYỄN THỊ LỤA (60 TUỔI — SỐ 12 ĐƯỜNG BỜ SÔNG)}
\begin{itemize}[leftmargin=1.5cm]
    \item \textbf{18:30 — 18:50:} Trời mưa to, xe máy Trần Ngọc Mai chở chồng là Lê Quang Vũ sang nhà Khang. Nghe tiếng Mai đập bàn quát tháo đòi đất đai di chúc. Đến 18:50 Mai nổ máy xe phóng về một mình dưới mưa, Vũ ở lại bên trong.
    \item \textbf{19:15:} Nghe tiếng Khang chửi bới rồi tiếng tát bốp bốp. Liền sau đó thấy Vũ ôm má trái, cổ áo xộc xệch hoảng loạn lao xe máy tháo chạy trong mưa.
    \item \textbf{20:00 — 20:15:} Lúc đang ngồi xem tập phim truyền hình trên kênh \texttt{TH3} (sóng truyền hình khu Bờ Sông phát sóng hoàn toàn bình thường), nghe tiếng cãi nhau lớn của đàn ông bên nhà Khang nhắc đến tên \textit{"thằng Huy"}, rồi tiếng \textit{"XOẢNG!"} vỡ vụn đồ gốm sứ. Đến 20:15 thấy một người đàn ông mặc áo bảo hộ công nhân xây dựng (Tùng) hớt hải chạy vội ra ngõ.
    \item \textbf{20:45:} Nhìn qua cửa sổ thấy một bóng người mặc áo gió trùm kín đầu màu xám đen đứng rình dưới gốc cây xoan đầu ngõ rồi lén lẻn vào nhà Khang.
    \item \textbf{06:30 — 06:45 ngày 25/07 (Phát hiện thi thể \& Báo án):} Sáng sớm ra quét ngõ, thấy cửa nhà Khang mở toang từ đêm qua, đèn phòng khách bật sáng, nước mưa tạt lênh láng. Bà Lụa sang gọi thì bàng hoàng phát hiện Khang nằm gục trên vũng máu bên cạnh bình trà vỡ, toàn thân co cứng. Bà hoảng hốt tri hô và gọi Công an Phường lúc 06:45.
\end{itemize}

\section*{II. LỜI KHAI ÔNG BÀ NGUYỄN VĂN TIẾN (BỐ MẸ TÙNG — SỐ 10 BỜ SÔNG)}
\begin{itemize}[leftmargin=1.5cm]
    \item \textbf{Quan hệ thời thơ ấu:} Ngày xưa thằng Khang hay chơi cùng hai đứa con trai nhà ông bà (Tùng và bé Gia Huy), đám trẻ thân nhau lắm. Từ sau biến cố đau lòng năm 1998 khi thằng út (Gia Huy) qua đời thì bọn trẻ không còn qua lại với nhau nữa.
    \item \textbf{Tính cách nạn nhân:} Khang từ nhỏ tính nết ngỗ ngược, hay gây gổ, lớn lên gia đình không còn giao thiệp.
    \item \textbf{Tình tiết gần đây:} Mấy hôm trước, gia đình có bảo con trai lớn (Tùng) tiện đường ghé qua nhà Khang gửi lời mời sang ăn đám giỗ em trai, nhưng không thấy Khang sang. Tùng hiện đi làm thợ xây dựng công trình, thỉnh thoảng mới về quê.
\end{itemize}

\vspace{0.8cm}

% --- CHỮ KÝ NGHỊ ĐỊNH 30 ---
\noindent
\begin{minipage}[t]{0.45\textwidth}
    \small
    \textbf{\textit{Nơi nhận:}}\\
    - Hồ sơ vụ án TEST-99;\\
    - Lưu: VT, CSHS.
\end{minipage}
\hfill
\begin{minipage}[t]{0.5\textwidth}
    \begin{center}
        \small
        \textbf{CÁN BỘ LẤY LỜI KHAI}\\
        \textit{(Ký, ghi rõ họ tên)}\\[1.5cm]
        \textbf{Trung úy Nguyễn Văn Hoàng}
    \end{center}
\end{minipage}

\end{document}
